% !TEX program = pdflatex
% arXiv submission source (English version), compiled with pdfLaTeX
\documentclass[11pt]{article}
\usepackage[a4paper,margin=2.5cm]{geometry}
\usepackage{amsmath,amssymb}
\usepackage{booktabs}
\usepackage{enumitem}
\usepackage{hyperref}
\usepackage{xcolor}
\newcommand{\infmark}{\textcolor{orange}{(\(\dagger\))}}

\title{Barriers, Work, and Returns:\\A Unified Socioeconomic Analytical Framework}
\author{Lei Dai}
\date{October 2026}

\begin{document}
\maketitle

\begin{abstract}
Starting from a fundamental question --- how do human markets come into existence --- this paper constructs a unified socioeconomic analytical framework. We treat trade as a state flow in which society moves from diversity toward sameness, define a barrier as an existing structure that maintains a potential differential against the homogenizing trend, and argue that a barrier holder's return is the compensation for its maintenance cost. The framework's main components are: (1) a three-dimensional description of barriers (breadth of domain, depth of subject-embeddedness, time structure of tolls) and a seven-stage sequence of attack, defense, and retreat along the embeddedness axis; (2) a unified ``work--crystallization'' theory that classifies individual capability and structural power as two phases of the same flow of work in two kinds of containers, yielding a four-variable representation of returns; and (3) boundary-layer propositions: homogenization is asymptotic rather than terminal, the lower bound of differentiation follows from the non-homogenizability of decision rights plus physical constraints, and automation is an endogenous, ever-lengthening river of externalizing work. On this basis we derive six strategy rules and test the framework against sixteen sets of documentary evidence and historical cases (hand-loom weavers, bank tellers and ATMs, radiology, nurses versus star doctors, $r>g$, among others); propositions for which no direct evidence yet exists and which rest on inference are honestly flagged with \infmark{} throughout the text. The paper's contribution is to unify three economic traditions --- division of labor, transaction costs, and rent --- on a single axis (subject-embeddedness), and to provide falsifiable, structured predictions for decision-making in the AI era.

\medskip
\noindent\textbf{Keywords:} transaction costs; barrier rent; work and crystallization; automation; labor markets in the AI era
\end{abstract}

\section{Introduction}
\subsection{Where the Question Starts}
How do markets for human trade come into being? The standard answer appeals to division of labor and self-interest: individuals produce different goods and profit by exchanging them. This paper adopts a more fundamental perspective: trade can be regarded as a state flow in which human society moves from diversity toward sameness; the mode of trade is infrastructure, and transaction costs are the cost of flowing across a given infrastructure --- the compensation for flattening differences. Historically, the barrier that acts as the main cost of flow keeps rotating: in the agrarian age it was the transport barrier; as technology advanced it became the information barrier; in the software era it became the production barrier. As generative AI pushes the cost of ``generation'' toward zero, the production barrier is collapsing --- what will be the most expensive thing next?

Pursuing this question to the end yields three fundamental questions: (1) How does the tension between differentiation and sameness drive social change? (2) What determines an individual's returns in this process? (3) How should an individual choose? This paper answers them in turn with the framework (Section~3), a theory of choice (Section~4), and empirical tests (Section~5).

\subsection{Relation to Economic Traditions}
The framework unifies three traditions: Smith's theory of the division of labor (the spatial distribution of differences), Coase's theory of transaction costs (barriers define the boundary of the market), and Ricardo's theory of rent (scarce barriers yield rent). Schumpeter's innovation rent, Polanyi's tacit knowledge, Bourdieu's conversion of forms of capital, and North's institutional theory then serve, respectively, as key inputs to the framework's dynamic layer, structural layer, unifying layer, and institutional boundary.

\subsection{Methodological Note}
This paper is a systematic attempt at theory construction: start from observation and propose a unifying hypothesis; converge the framework through sustained dialogue with historical and contemporary cases; then derive strategies and submit them to stress-testing by counterexamples. Every empirical claim is given a documentary or data source wherever possible; propositions for which no direct evidence can be found are flagged with \infmark{} and their nature is clarified collectively in Section~6.

\section{Concepts and Definitions}
\begin{description}[leftmargin=2.2em]
\item[Potential differential] The differential distribution of the world along dimensions such as goods, information, capability, and trust. Differentials drive flows (trade, learning, imitation) and are the prime mover of social change.
\item[Barrier] An existing structure that maintains a potential differential against the homogenizing flow. The barrier holder pays a maintenance cost, and its return is the compensation for that cost (rent).
\item[Work] The sustained input of labor, attention, and energy --- the ``currency'' of all barriers.
\item[Crystallization] The conversion of flowing work, through a registration system (property rights, diplomas, brands, platform weights, etc.), into a stock that can be stored, transferred, and inherited.
\item[Container] The place where a crystallization is registered: an individual body, an organization, or an institutional position; container ownership determines the security and inheritability of the stock.
\item[Depreciation] The decay of a stock through forgetting, institutional change, and network migration; maintenance-type work hedges against depreciation.
\end{description}

\section{The Framework}
\subsection{Ontological Layer: Differentials and the Prime Mover}
Differentials exist in the world, and differentials drive flows. Homogenization is the prime mover of social change, but it does not produce monotonic convergence: barriers maintain differentials against the flow, so there is persistent pressure on both sides of every barrier, and society pays two kinds of cost at the barrier --- the cost of crossing (transaction costs) and the cost of maintenance. The barrier holder's return is the compensation for the maintenance cost.

\subsection{Structural Layer: Three Dimensions and the Attack--Retreat Sequence}
The complete description of a barrier is a product of three dimensions: breadth of domain $\times$ depth of subject-embeddedness $\times$ time structure of tolls. Attack, defense, and retreat form a continuous sequence along the embeddedness axis:
\[
\text{Goods} \to \text{Information} \to \text{Knowledge} \to \text{Trust} \to \text{Interpretation rights} \to \text{Governance rights} \to \text{Decision rights}
\]
The distinction between adjacent levels reduces to a single axis: the depth to which the barrier is embedded in the subject (detachability decreasing), whose side facets include decreasing replicability, decreasing verifiability, increasing consensus requirements, increasing replacement cost, and a decreasing capacity to ``collect rent while lying down.'' Homogenization's tactic against relational barriers is to objectify first (trust $\to$ scores, law $\to$ code) and homogenize second; retreat therefore keeps advancing toward meta-levels.

\subsection{Dynamic Layer: Cost, Competition, and Infrastructure}
The rent-maximizing form of a barrier is a controlled leak: fully impermeable and the flow disappears; fully permeable and the differential disappears; the rent-maximizing calibration point lies in between. Infrastructure is a composite of ``toll collection $+$ wall-demolition services''; competition among barriers of the same kind is decided by net barrier cost --- the winner's rent minus the value of its wall-demolition service to others must be smaller than the network's cost of building an alternative\infmark. When domains do not fully overlap, old barriers survive on residual dimensions (cash in anonymous scenarios of the e-payment era). The replacement channel of institutional barriers is political, with lags of up to a century --- the framework has directional but not temporal predictive power at the institutional layer.

\subsection{Unifying Layer: The Work--Crystallization Theory}
The two great classes of barriers (individual capability and structural power) are unified as phases of work: gaseous (present only while the work is being done, constrained by 24 hours), liquid (the gray zone: organizational culture, deep relationships), and solid (storable, transferable, inheritable). Crystallization is performed by registration systems; the critical parameter is container ownership --- the crystallization of strong power sits in self-owned registration systems, while the crystallization of weak power (e.g., platform influencers) sits in rented containers and can be confiscated in an instant. The solid phase exempts one from being present, but not from working: work shifts from generative to maintenance-type, and depreciation comes from three directions --- forgetting, institutional change, and network migration.

The master formula\infmark:
\[
\text{Returns} = f\big(\ \underbrace{\text{Historical work input} \times \text{Crystallization efficiency} \times \text{Container ownership}}_{\text{formation of the stock}},\ \underbrace{\text{Maintenance work} \div \text{Depreciation rate}}_{\text{persistence of the stock}}\ \big)
\]
Equivalent reading: returns $\approx$ degree of dependence of the system on you $\times$ the speed at which that dependence can be dissolved $\times$ pricing power $-$ maintenance costs. The four variables cover all attribute differences between the two classes: Class 1 (individual capability) has the body as its container --- small capacity, fast depreciation, non-transferable, returns linear in 24 hours; Class 2 (structural power) has the structure as its container, whose maintenance cost can be externalized to the system --- the only form of returns that can exceed individual service limits without bound.

\subsection{Boundary Layer: Ripples, the Lower Bound, and Automation}
Ripples $=$ Gradient $\times$ Search. The gradient (fuel) comes entirely from the Earth--Sun physical layer; the search (ignition) is implemented by materialized information systems --- brains, scientific communities, R\&D organizations --- with consciousness as the searcher rather than the prime mover. In a homogeneous system, ripples are endogenous: the homogenizing offensive pushes the system away from equilibrium, and tiny fluctuations are spontaneously amplified into new barriers (spontaneous symmetry breaking) --- flattening and creating are two phases of the same process\infmark.

The lower bound of differentiation is above zero: homogenization is asymptotic, never terminal. Its strict form is the non-homogenizability of decision rights (one decision cannot be made by everyone simultaneously) \cite{arrow1951,radner1993}, plus physical constraints (embodiment, energy gradients). The lower bound is not a ruin; it is the last position each generation stands on.

Automation is the extracorporeal crystallization of work (an endogenous tendency): memory $\to$ writing $\to$ databases, computation $\to$ computers, generation $\to$ models. With each externalization, the price of that capability falls to zero, and the human position retreats one level along the embeddedness axis. The executor of each automation wave is systematically the target of the next (blacksmith $\to$ mechanic $\to$ programmer) --- zero counterexamples in two centuries\infmark. The AI era's specificity is only a change of parameters: depreciation is an order of magnitude faster, and externalization touches ``generation'' for the first time; but generation was never the lower bound of barriers --- only the latest position of the formalizable frontier.

\section{A Theory of Choice}
\subsection{The Degrees of Freedom of Choice}
Choice is effective on only two degrees of freedom: position (the three-dimensional coordinates of a barrier) and search (the betting rate on new ripples). Choice determines the structure and exposure of one's returns, not the overall water level. Embeddedness and binding are two readings of the same variable: any position with persistently high returns is in fact paid for with freedom --- the high returns at a barrier are, in essence, society's compensation for ``being willing to be bound there, working continuously.''

\subsection{Six Core Strategies}
\begin{enumerate}[leftmargin=2em]
\item Position on a layer with ``high pressure $\times$ slow attack speed'': migrate toward layers that are ``hard to learn, hard to imitate, hard to take away'' (trust, judgment, embodiment, tacit knowledge, responsibility-bearing signatures).
\item Flowing with the current beats blocking it: make things that make other people's flows cheaper, and collect tolls on the channel; do not build roadblocks.
\item The objectification test: if the core of a barrier can be fully described by a database, a scoring system, or a prompt, its rent is already decaying.
\item Convert one-time income into a toll gate: restructure one-off income into a checkpoint that skims every flow.
\item Search strategy --- ``rather be wrong early than right late'': the first fluctuation to be amplified takes the positive feedback and path dependence; the expectation depends on who gets amplified first, not on who ends up best. Bet small, bet often, bet broadly.
\item Manage a portfolio across embeddedness: homogenizable accounts (for liquidity) and non-homogenizable accounts (holding durable barriers) allocated along the life cycle.
\end{enumerate}

\subsection{Operationalizing the Four Variables}
Three questions about an occupation: in whose name is your work crystallized; the ratio of flow to stock in your pay (what can you take with you if you leave in three years); which depreciation bracket your skill sits in and how generatable it is. Reviewing employers and platforms: crystallization clauses (authorship, client ownership, non-compete boundaries); container ownership (multi-platform distribution $=$ opening multiple passbooks). Four principles of asset allocation: continuously convert flow income into solid stocks; diversify containers; match depreciation rates to age; always retain gaseous capability. The minimalist daily strategy in the AI era: leave generation to the machine, keep judgment to yourself, sign your own name on responsibility.

\subsection{Validity Across Eras}
The six strategies are structurally identical across the four automation waves of agriculture, physical labor, computation, and information; the wrong strategies are also identical (deepening into an obsolescent configuration layer, keeping one's passbook in the employer's container, blocking trends with a barrier). Arkwright in 1769 (production to the machine, judgment and responsibility to himself, building his own container) was the earliest complete implementation of this strategy structure; the Luddite movement and the guild path produced zero successful cases in two centuries.

\section{Empirical Evidence and Case Tests}
Table~\ref{tab:evidence} summarizes the representative evidence for each proposition (the full mapping is in the online appendix). Key data points: hand-loom weavers' weekly wage was 240 pence in 1806, falling to 99 pence by 1820 --- more than 25\% below factory workers (Acemoglu, 2024, w32416); ATMs reduced tellers per branch from about 20 to 13 while total employment did not fall (Bessen, 2015); U.S. radiologist employment grew 17.3\% between 2014 and 2023; the employment gap for 22--25-year-olds in the most AI-exposed occupations reached 19\% (Canaries, 2026); U.S. agricultural employment fell from about 40\% of total employment in 1900 to about 2\% in 2000.

\begin{table}[h]
\centering\small
\caption{Proposition--evidence mapping (abridged)}
\label{tab:evidence}
\begin{tabular}{p{0.42\textwidth}p{0.48\textwidth}}
\toprule
Proposition & Representative evidence (strength)\\
\midrule
Trade and division of labor are the exchange and flattening of differences & Smith (1776); Coase (1937, 1960); Ricardo (1817) (classic)\\
Transport-barrier removal and barrier rotation & Levinson (2006) (data)\\
Codifiable knowledge depreciates; tacit knowledge is protected & Polanyi (1958, 1966); Canaries (2026) (classic/empirical)\\
High rent breeds attack (counter-current ripples) & Allen; Jevons (1865) (empirical/classic)\\
Barrier costs fall, total flow surges & Bessen (2015); radiology employment data (empirical/data)\\
Political lag of institutional barriers & North \& Weingast (1989) (empirical, contested)\\
Capital is a relation, not a thing & Bourdieu (1986); Marx (classic)\\
Decision rights cannot be homogenized & Arrow (1951); Radner (1993) (classic)\\
Path dependence and break points & David (1985); Arthur (1989) (empirical, contested)\\
Responsibility and judgment survive the waves & Hinton 2016 and HBR 2026 retrospective; SOX (2002) (report/institution)\\
Network effects amplify back; AI sweeps only Class 1 & Katz \& Shapiro (1985); Canaries (2026) (classic/data)\\
The unit of collapse is the barrier component & Autor et al.\ (2003); Acemoglu \& Restrepo (2018, 2022) (empirical)\\
Monopsony intercepts rent & Manning (2003) (classic)\\
$r>g$ & Piketty (2014); Piketty \& Zucman (2014) (empirical)\\
Strategies are invariant across eras & USDA agricultural employment data; Brynjolfsson et al.\ (2025) (data/empirical)\\
Thermodynamic analogies & Prigogine (1977); Georgescu-Roegen (1971); Hawking (1974) (classic)\\
\bottomrule
\end{tabular}
\end{table}

Digesting counterexamples is the framework's main mechanism of upgrading. The nurse/teacher cases exposed a missing ``pricing power'' factor in the formula (ownership of the valve); the ``capital idles and wins'' case exposed the hidden assumption that ``maintenance work can be externalized''; both were digested inside the formula, which after correction reads: long-run returns $\approx$ societal dependence $\times$ time $\times$ valve $-$ the bill.

\section{Scope Limits and Honest Flags}
The framework's known limits: (1) break points are unpredictable --- which fluctuation gets spontaneously amplified is near a lottery; the formula describes the persistence of rent after the break, not the break itself; (2) temporal prediction fails at the institutional layer --- the lag is decided by politics; (3) the depreciation of a rented container contains the container owner's will as a variable and cannot be modeled separately.

The following propositions currently lack direct empirical support and are internal inferences or synthetic innovations of the framework\infmark: the net-barrier-cost criterion for competition among registration systems; the unified theory of registration systems as crystallizers of work, and the four-variable master formula; the quantitative relation between maintenance work and the depreciation rate; the recursion ``executor of an automation wave $\to$ target of the next''; the prediction that ``the barriers remaining in the AI era $=$ verification $+$ responsibility $+$ embodiment $+$ trust''; and the strong formulation that ``flattening and creating are the same phase'' of ripples. These propositions constitute the framework's falsifiable frontier.

\section{Testable Assertions}
To test the framework's predictive power, we convert its core extrapolative propositions into six falsifiable assertions. The assertions follow the public conventions of forecasting platforms (the criteria--deadline--edge-cases triad of the Metaculus question-writing guide\footnote{\url{https://metaculus-metaculus.mintlify.app/guides/question-types}}; the ``confidence-tagged assertion $+$ next-year retrospective'' format of the Astral Codex Ten annual prediction contest\footnote{\url{https://www.astralcodexten.com/p/2023-prediction-contest}; a sample retrospective at \url{https://www.astralcodexten.com/p/grading-my-2018-predictions-for-2023}}). Each assertion is bound to a third-party data source with a clear release cycle; deadlines are back-derived from data availability. We recommend calibrating with Brier scores year by year --- if the calibration curve deviates systematically, the framework should be revised, not the assertions.

\begin{description}[leftmargin=2.2em,style=nextline]
\item[Assertion 1 (Widening youth employment gap in codified knowledge, confidence 60\%)]
The ``keeping-pace gap'' in employment for 22--25-year-olds in the most AI-exposed occupations will widen from 19\% in the June 2026 data period to 25\% or more. Win: an updated version of the Canaries study (or a same-specification replication on ADP data) reports a gap $\geq 25\%$. Lose: the gap falls back below 19\%, or stays flat for two consecutive years, or the tacit-knowledge occupation group shows a gap of the same magnitude. Deadline 2028-06-30 (Canaries is revised every August; this date covers both the 2027.8 and 2028 updates). Data source: Stanford Digital Economy Lab, \textit{Canaries in the Coal Mine}.
\item[Assertion 2 (Barrier components turn ``expensive'': the radiology pattern repeats, confidence 70\%)]
After AI penetrates the codifiable sub-tasks of an occupation, the occupation's total demand and price rise rather than fall: for U.S. radiologists (SOC 29-1224), at least one of the OEWS employment count and the median annual wage grows faster than ``physicians overall.'' Win: at least one of the two indicators grows no slower than physicians overall. Lose: both indicators grow more slowly than physicians overall, or radiology employment falls in absolute terms. Deadline 2028-06-30 (OEWS is annual, with a May reference period and a $\sim$12-month publication lag; this date incorporates the new data for the May 2027 reference period). Data source: BLS OEWS.
\item[Assertion 3 (AI sweeps only Class-1 barriers, confidence 60\%)]
Embodied/service occupations (nursing, repair, personal care) see average hourly wage growth at least equal to the all-industry average, with no large-scale, AI-attributable displacement. Win: in the CES monthly industry wage series, cumulative growth of health care \& social assistance $\geq$ the private nonfarm total. Lose: growth persistently more than 2 percentage points below the total, or a visible contraction of nursing/repair employment. Deadline 2027-12-31 (CES is released monthly with a $\sim$1-month lag; a one-year window suffices to judge direction). Data source: BLS CES.
\item[Assertion 4 (The hollowing of the middle deepens, confidence 65\%)]
The 90/50 percentile ratio of the U.S. wage distribution keeps rising, and the share of mid-skill routine jobs keeps falling. Win: CPS ASEC or OEWS quantile data show the 90/50 ratio rising for two consecutive periods. Lose: two consecutive periods of decline, or a rebound in the mid-skill share. Deadline 2028-06-30 (``two consecutive periods'' requires ASEC 2026 data --- released September 2027 --- and the OEWS May 2027 reference period --- released May 2028 --- two annual observations). Data source: Census CPS ASEC; BLS OEWS.
\item[Assertion 5 (The next target of the recursion appears, confidence 50\%)]
Today's high-paid configuration-layer jobs created by AI (prompt engineering, AI content operations, junior AI operations) will be absorbed by built-in tool capabilities: postings for such jobs on major hiring platforms fall by 50\% or more from their 2025--2026 peak. Win: the Indeed Hiring Lab job-posting index (or LinkedIn data on the same specification) shows a decline $\geq 50\%$. Lose: postings rise instead, or the jobs upgrade into more senior roles with total volume at a new high. Deadline 2028-06-30 (``absorption by built-in tools'' takes roughly two product cycles). This assertion tests an inference-level proposition \infmark{} and has the weakest data foundation; its confidence is correspondingly the lowest. Data source: Indeed Hiring Lab.
\item[Assertion 6 (The political response on container ownership, confidence 50\%)]
At least one major economy (US/EU/China) brings into force a new statute targeting AI platforms or large platforms, requiring any one of: creator data portability, algorithmic-weight auditability, or traceability of model outputs. Win: a publicly verifiable enacted statute (checkable on EUR-Lex or official national legal databases). Lose: no such statute takes effect. Deadline 2028-06-30 (legislative cycles run in years; this assertion bets on direction, not on timing --- the two-year window corresponds to the lower bound of the maximum political lag the framework allows). The winning criterion deliberately excludes portability clauses already satisfied by GDPR/DMA, counting only new rules in the AI-platform direction.
\end{description}

\section{Conclusion}
Market trade is the flow of potential differentials; a barrier is the crystallization against the current; rent is the compensation for maintaining the crystallization; work is the currency of all barriers; the registration system is the exchange; depreciation is entropy; automation is the ever-lengthening river of externalization. The four variables of individual fate: total historical work input, crystallization efficiency, container ownership, and the ratio of maintenance work to the depreciation rate. To the original question --- will returns that exceed personal limits eventually be eliminated? --- the corrected answer is: every return that does not attach to continuing personal input will be eliminated, and every return that does attach to it already is, by itself, personal service.

\begin{thebibliography}{99}\small
\bibitem{smith1776} Smith, A. (1776). \textit{An Inquiry into the Nature and Causes of the Wealth of Nations}.
\bibitem{coase1937} Coase, R. H. (1937). The Nature of the Firm. \textit{Economica}, 4(16), 386--405.
\bibitem{coase1960} Coase, R. H. (1960). The Problem of Social Cost. \textit{Journal of Law and Economics}, 3, 1--44.
\bibitem{ricardo1817} Ricardo, D. (1817). \textit{On the Principles of Political Economy and Taxation}.
\bibitem{schumpeter1942} Schumpeter, J. A. (1942). \textit{Capitalism, Socialism and Democracy}.
\bibitem{polanyi1958} Polanyi, M. (1958). \textit{Personal Knowledge}. University of Chicago Press.
\bibitem{polanyi1966} Polanyi, M. (1966). \textit{The Tacit Dimension}. Doubleday.
\bibitem{arrow1951} Arrow, K. J. (1951). \textit{Social Choice and Individual Values}. Wiley.
\bibitem{radner1993} Radner, R. (1993). The Organization of Decentralized Information Processing. \textit{Econometrica}, 61(5), 1109--1146.
\bibitem{marschak1972} Marschak, J., \& Radner, R. (1972). \textit{Economic Theory of Teams}. Yale University Press.
\bibitem{almt2003} Autor, D. H., Levy, F., \& Murnane, R. J. (2003). The Skill Content of Recent Technological Change. \textit{Quarterly Journal of Economics}, 118(4), 1279--1333.
\bibitem{autor2015} Autor, D. H. (2015). Why Are There Still So Many Jobs? \textit{Journal of Economic Perspectives}, 29(3), 3--30.
\bibitem{ar2018} Acemoglu, D., \& Restrepo, P. (2018). The Race between Man and Machine. \textit{American Economic Review}, 108(6), 1488--1542.
\bibitem{ar2022} Acemoglu, D., \& Restrepo, P. (2022). Tasks, Automation, and the Rise in US Wage Inequality. \textit{Econometrica}, 90(5), 1973--2016.
\bibitem{acemoglu2024} Acemoglu, D. (2024). The Simple Macroeconomics of AI. NBER Working Paper No. 32416.
\bibitem{allen} Allen, R. C. The Hand-Loom Weaver and the Power Loom. Working paper.
\bibitem{bessen2015} Bessen, J. (2015). \textit{Learning by Doing: The Real Connection between Innovation, Wages, and Wealth}. Yale University Press.
\bibitem{levinson2006} Levinson, M. (2006). \textit{The Box: How the Shipping Container Made the World Smaller and the World Economy Bigger}. Princeton University Press.
\bibitem{hinton2016} Hinton, G. (2016). Public remarks: ``People should stop training radiologists now.''; retrospective in HBR (2026).
\bibitem{sox2002} Sarbanes--Oxley Act (2002). \S 302/906.
\bibitem{canaries2026} Brynjolfsson, E., Chandar, B., \& Chen, L. (2026). Canaries in the Coal Mine? Stanford Digital Economy Lab (August 2026 revision).
\bibitem{nw1989} North, D. C., \& Weingast, B. R. (1989). Constitutions and Commitment. \textit{Journal of Economic History}, 49(4), 803--832.
\bibitem{bourdieu1986} Bourdieu, P. (1986). The Forms of Capital. In \textit{Handbook of Theory and Research for the Sociology of Education}.
\bibitem{manning2003} Manning, A. (2003). \textit{Monopsony in Motion}. Cambridge University Press.
\bibitem{piketty2014} Piketty, T. (2014). \textit{Capital in the Twenty-First Century}. Harvard University Press.
\bibitem{pz2014} Piketty, T., \& Zucman, G. (2014). Capital is Back. \textit{Quarterly Journal of Economics}, 129(3), 1255--1310.
\bibitem{david1985} David, P. A. (1985). Clio and the Economics of QWERTY. \textit{AER P\&P}, 75(2), 332--337.
\bibitem{arthur1989} Arthur, W. B. (1989). Competing Technologies, Increasing Returns, and Lock-In by Historical Events. \textit{Economic Journal}, 99, 116--131.
\bibitem{ks1985} Katz, M. L., \& Shapiro, C. (1985). Network Externalities, Competition, and Compatibility. \textit{American Economic Review}, 75(3), 424--440.
\bibitem{blr2025} Brynjolfsson, E., Li, D., \& Raymond, L. (2025). Generative AI at Work. \textit{Quarterly Journal of Economics}, 140(2), 889--942.
\bibitem{jevons1865} Jevons, W. S. (1865). \textit{The Coal Question}.
\bibitem{prigogine1977} Prigogine, I. (1977). Nobel Prize Lecture (dissipative structures).
\bibitem{georgescu1971} Georgescu-Roegen, N. (1971). \textit{The Entropy Law and the Economic Process}. Harvard University Press.
\bibitem{hawking1974} Hawking, S. W. (1974). Black Hole Explosions? \textit{Nature}, 248, 30--31.
\bibitem{ar2012} Acemoglu, D., \& Robinson, J. A. (2012). \textit{Why Nations Fail}. Crown.
\end{thebibliography}

\end{document}
